% Options for packages loaded elsewhere
\PassOptionsToPackage{unicode}{hyperref}
\PassOptionsToPackage{hyphens}{url}
%
\documentclass[
]{article}
\usepackage{amsmath,amssymb}
\usepackage{iftex}
\ifPDFTeX
  \usepackage[T1]{fontenc}
  \usepackage[utf8]{inputenc}
  \usepackage{textcomp} % provide euro and other symbols
\else % if luatex or xetex
  \usepackage{unicode-math} % this also loads fontspec
  \defaultfontfeatures{Scale=MatchLowercase}
  \defaultfontfeatures[\rmfamily]{Ligatures=TeX,Scale=1}
\fi
\usepackage{lmodern}
\ifPDFTeX\else
  % xetex/luatex font selection
\fi
% Use upquote if available, for straight quotes in verbatim environments
\IfFileExists{upquote.sty}{\usepackage{upquote}}{}
\IfFileExists{microtype.sty}{% use microtype if available
  \usepackage[]{microtype}
  \UseMicrotypeSet[protrusion]{basicmath} % disable protrusion for tt fonts
}{}
\makeatletter
\@ifundefined{KOMAClassName}{% if non-KOMA class
  \IfFileExists{parskip.sty}{%
    \usepackage{parskip}
  }{% else
    \setlength{\parindent}{0pt}
    \setlength{\parskip}{6pt plus 2pt minus 1pt}}
}{% if KOMA class
  \KOMAoptions{parskip=half}}
\makeatother
\usepackage{xcolor}
\usepackage[margin=1in]{geometry}
\usepackage{color}
\usepackage{fancyvrb}
\newcommand{\VerbBar}{|}
\newcommand{\VERB}{\Verb[commandchars=\\\{\}]}
\DefineVerbatimEnvironment{Highlighting}{Verbatim}{commandchars=\\\{\}}
% Add ',fontsize=\small' for more characters per line
\usepackage{framed}
\definecolor{shadecolor}{RGB}{248,248,248}
\newenvironment{Shaded}{\begin{snugshade}}{\end{snugshade}}
\newcommand{\AlertTok}[1]{\textcolor[rgb]{0.94,0.16,0.16}{#1}}
\newcommand{\AnnotationTok}[1]{\textcolor[rgb]{0.56,0.35,0.01}{\textbf{\textit{#1}}}}
\newcommand{\AttributeTok}[1]{\textcolor[rgb]{0.13,0.29,0.53}{#1}}
\newcommand{\BaseNTok}[1]{\textcolor[rgb]{0.00,0.00,0.81}{#1}}
\newcommand{\BuiltInTok}[1]{#1}
\newcommand{\CharTok}[1]{\textcolor[rgb]{0.31,0.60,0.02}{#1}}
\newcommand{\CommentTok}[1]{\textcolor[rgb]{0.56,0.35,0.01}{\textit{#1}}}
\newcommand{\CommentVarTok}[1]{\textcolor[rgb]{0.56,0.35,0.01}{\textbf{\textit{#1}}}}
\newcommand{\ConstantTok}[1]{\textcolor[rgb]{0.56,0.35,0.01}{#1}}
\newcommand{\ControlFlowTok}[1]{\textcolor[rgb]{0.13,0.29,0.53}{\textbf{#1}}}
\newcommand{\DataTypeTok}[1]{\textcolor[rgb]{0.13,0.29,0.53}{#1}}
\newcommand{\DecValTok}[1]{\textcolor[rgb]{0.00,0.00,0.81}{#1}}
\newcommand{\DocumentationTok}[1]{\textcolor[rgb]{0.56,0.35,0.01}{\textbf{\textit{#1}}}}
\newcommand{\ErrorTok}[1]{\textcolor[rgb]{0.64,0.00,0.00}{\textbf{#1}}}
\newcommand{\ExtensionTok}[1]{#1}
\newcommand{\FloatTok}[1]{\textcolor[rgb]{0.00,0.00,0.81}{#1}}
\newcommand{\FunctionTok}[1]{\textcolor[rgb]{0.13,0.29,0.53}{\textbf{#1}}}
\newcommand{\ImportTok}[1]{#1}
\newcommand{\InformationTok}[1]{\textcolor[rgb]{0.56,0.35,0.01}{\textbf{\textit{#1}}}}
\newcommand{\KeywordTok}[1]{\textcolor[rgb]{0.13,0.29,0.53}{\textbf{#1}}}
\newcommand{\NormalTok}[1]{#1}
\newcommand{\OperatorTok}[1]{\textcolor[rgb]{0.81,0.36,0.00}{\textbf{#1}}}
\newcommand{\OtherTok}[1]{\textcolor[rgb]{0.56,0.35,0.01}{#1}}
\newcommand{\PreprocessorTok}[1]{\textcolor[rgb]{0.56,0.35,0.01}{\textit{#1}}}
\newcommand{\RegionMarkerTok}[1]{#1}
\newcommand{\SpecialCharTok}[1]{\textcolor[rgb]{0.81,0.36,0.00}{\textbf{#1}}}
\newcommand{\SpecialStringTok}[1]{\textcolor[rgb]{0.31,0.60,0.02}{#1}}
\newcommand{\StringTok}[1]{\textcolor[rgb]{0.31,0.60,0.02}{#1}}
\newcommand{\VariableTok}[1]{\textcolor[rgb]{0.00,0.00,0.00}{#1}}
\newcommand{\VerbatimStringTok}[1]{\textcolor[rgb]{0.31,0.60,0.02}{#1}}
\newcommand{\WarningTok}[1]{\textcolor[rgb]{0.56,0.35,0.01}{\textbf{\textit{#1}}}}
\usepackage{graphicx}
\makeatletter
\newsavebox\pandoc@box
\newcommand*\pandocbounded[1]{% scales image to fit in text height/width
  \sbox\pandoc@box{#1}%
  \Gscale@div\@tempa{\textheight}{\dimexpr\ht\pandoc@box+\dp\pandoc@box\relax}%
  \Gscale@div\@tempb{\linewidth}{\wd\pandoc@box}%
  \ifdim\@tempb\p@<\@tempa\p@\let\@tempa\@tempb\fi% select the smaller of both
  \ifdim\@tempa\p@<\p@\scalebox{\@tempa}{\usebox\pandoc@box}%
  \else\usebox{\pandoc@box}%
  \fi%
}
% Set default figure placement to htbp
\def\fps@figure{htbp}
\makeatother
\setlength{\emergencystretch}{3em} % prevent overfull lines
\providecommand{\tightlist}{%
  \setlength{\itemsep}{0pt}\setlength{\parskip}{0pt}}
\setcounter{secnumdepth}{-\maxdimen} % remove section numbering
\usepackage{graphicx}
\usepackage{float}
\usepackage{tikz}
\usepackage{geometry}
\usepackage{bookmark}
\IfFileExists{xurl.sty}{\usepackage{xurl}}{} % add URL line breaks if available
\urlstyle{same}
\hypersetup{
  hidelinks,
  pdfcreator={LaTeX via pandoc}}

\author{}
\date{\vspace{-2.5em}}

\begin{document}

\begin{Shaded}
\begin{Highlighting}[]
\FunctionTok{options}\NormalTok{(}\AttributeTok{OutDec =} \StringTok{"."}\NormalTok{)  }\CommentTok{\# Asegurar punto decimal en este chunk}

\CommentTok{\# Establecer semilla aleatoria para reproducibilidad}
\CommentTok{\#set.seed(sample(1:10000, 1))}

\CommentTok{\# Aleatorización del contexto geométrico}
\NormalTok{terminos\_triangulo }\OtherTok{\textless{}{-}} \FunctionTok{c}\NormalTok{(}\StringTok{"triángulo"}\NormalTok{, }\StringTok{"polígono triangular"}\NormalTok{, }\StringTok{"figura triangular"}\NormalTok{, }\StringTok{"forma triangular"}\NormalTok{)}
\NormalTok{termino\_triangulo }\OtherTok{\textless{}{-}} \FunctionTok{sample}\NormalTok{(terminos\_triangulo, }\DecValTok{1}\NormalTok{)}

\NormalTok{terminos\_altura }\OtherTok{\textless{}{-}} \FunctionTok{c}\NormalTok{(}\StringTok{"alturas"}\NormalTok{, }\StringTok{"altitudes"}\NormalTok{, }\StringTok{"perpendiculares"}\NormalTok{, }\StringTok{"líneas perpendiculares"}\NormalTok{)}
\NormalTok{termino\_altura }\OtherTok{\textless{}{-}} \FunctionTok{sample}\NormalTok{(terminos\_altura, }\DecValTok{1}\NormalTok{)}

\NormalTok{terminos\_punto }\OtherTok{\textless{}{-}} \FunctionTok{c}\NormalTok{(}\StringTok{"punto"}\NormalTok{, }\StringTok{"lugar"}\NormalTok{, }\StringTok{"posición"}\NormalTok{, }\StringTok{"ubicación"}\NormalTok{)}
\NormalTok{termino\_punto }\OtherTok{\textless{}{-}} \FunctionTok{sample}\NormalTok{(terminos\_punto, }\DecValTok{1}\NormalTok{)}

\NormalTok{terminos\_geometrico }\OtherTok{\textless{}{-}} \FunctionTok{c}\NormalTok{(}\StringTok{"geométrico"}\NormalTok{, }\StringTok{"espacial"}\NormalTok{, }\StringTok{"bidimensional"}\NormalTok{, }\StringTok{"plano"}\NormalTok{)}
\NormalTok{termino\_geometrico }\OtherTok{\textless{}{-}} \FunctionTok{sample}\NormalTok{(terminos\_geometrico, }\DecValTok{1}\NormalTok{)}

\NormalTok{terminos\_cruzan }\OtherTok{\textless{}{-}} \FunctionTok{c}\NormalTok{(}\StringTok{"cruzan"}\NormalTok{, }\StringTok{"intersectan"}\NormalTok{, }\StringTok{"encuentran"}\NormalTok{, }\StringTok{"convergen"}\NormalTok{)}
\NormalTok{termino\_cruzan }\OtherTok{\textless{}{-}} \FunctionTok{sample}\NormalTok{(terminos\_cruzan, }\DecValTok{1}\NormalTok{)}

\NormalTok{terminos\_dibujado }\OtherTok{\textless{}{-}} \FunctionTok{c}\NormalTok{(}\StringTok{"dibujado"}\NormalTok{, }\StringTok{"trazado"}\NormalTok{, }\StringTok{"construido"}\NormalTok{, }\StringTok{"representado"}\NormalTok{)}
\NormalTok{termino\_dibujado }\OtherTok{\textless{}{-}} \FunctionTok{sample}\NormalTok{(terminos\_dibujado, }\DecValTok{1}\NormalTok{)}

\CommentTok{\# Aleatorización de letras para los vértices del triángulo}
\NormalTok{conjuntos\_vertices }\OtherTok{\textless{}{-}} \FunctionTok{list}\NormalTok{(}
  \FunctionTok{c}\NormalTok{(}\StringTok{"M"}\NormalTok{, }\StringTok{"O"}\NormalTok{, }\StringTok{"P"}\NormalTok{),  }\CommentTok{\# Conjunto original}
  \FunctionTok{c}\NormalTok{(}\StringTok{"A"}\NormalTok{, }\StringTok{"B"}\NormalTok{, }\StringTok{"C"}\NormalTok{),}
  \FunctionTok{c}\NormalTok{(}\StringTok{"D"}\NormalTok{, }\StringTok{"E"}\NormalTok{, }\StringTok{"F"}\NormalTok{),}
  \FunctionTok{c}\NormalTok{(}\StringTok{"R"}\NormalTok{, }\StringTok{"S"}\NormalTok{, }\StringTok{"T"}\NormalTok{),}
  \FunctionTok{c}\NormalTok{(}\StringTok{"X"}\NormalTok{, }\StringTok{"Y"}\NormalTok{, }\StringTok{"Z"}\NormalTok{),}
  \FunctionTok{c}\NormalTok{(}\StringTok{"L"}\NormalTok{, }\StringTok{"N"}\NormalTok{, }\StringTok{"Q"}\NormalTok{),}
  \FunctionTok{c}\NormalTok{(}\StringTok{"U"}\NormalTok{, }\StringTok{"V"}\NormalTok{, }\StringTok{"W"}\NormalTok{),}
  \FunctionTok{c}\NormalTok{(}\StringTok{"G"}\NormalTok{, }\StringTok{"H"}\NormalTok{, }\StringTok{"I"}\NormalTok{)}
\NormalTok{)}

\NormalTok{vertices\_seleccionados }\OtherTok{\textless{}{-}} \FunctionTok{sample}\NormalTok{(conjuntos\_vertices, }\DecValTok{1}\NormalTok{)[[}\DecValTok{1}\NormalTok{]]}
\NormalTok{vertice1 }\OtherTok{\textless{}{-}}\NormalTok{ vertices\_seleccionados[}\DecValTok{1}\NormalTok{]  }\CommentTok{\# Vértice izquierdo}
\NormalTok{vertice2 }\OtherTok{\textless{}{-}}\NormalTok{ vertices\_seleccionados[}\DecValTok{2}\NormalTok{]  }\CommentTok{\# Vértice superior (ortocentro)}
\NormalTok{vertice3 }\OtherTok{\textless{}{-}}\NormalTok{ vertices\_seleccionados[}\DecValTok{3}\NormalTok{]  }\CommentTok{\# Vértice derecho}

\CommentTok{\# Aleatorización de la letra para el pie de altura}
\NormalTok{letras\_disponibles }\OtherTok{\textless{}{-}} \FunctionTok{setdiff}\NormalTok{(LETTERS, vertices\_seleccionados)}
\NormalTok{pie\_altura }\OtherTok{\textless{}{-}} \FunctionTok{sample}\NormalTok{(letras\_disponibles, }\DecValTok{1}\NormalTok{)}

\CommentTok{\# Aleatorización de parámetros geométricos para el triángulo}
\CommentTok{\# Coordenadas del triángulo (manteniendo proporciones razonables)}
\NormalTok{coord\_base }\OtherTok{\textless{}{-}} \FunctionTok{sample}\NormalTok{(}\FunctionTok{c}\NormalTok{(}\DecValTok{4}\NormalTok{, }\DecValTok{5}\NormalTok{, }\DecValTok{6}\NormalTok{, }\DecValTok{7}\NormalTok{, }\DecValTok{8}\NormalTok{), }\DecValTok{1}\NormalTok{)  }\CommentTok{\# Ancho de la base}
\NormalTok{altura\_triangulo }\OtherTok{\textless{}{-}} \FunctionTok{sample}\NormalTok{(}\FunctionTok{c}\NormalTok{(}\DecValTok{3}\NormalTok{, }\DecValTok{4}\NormalTok{, }\DecValTok{5}\NormalTok{, }\DecValTok{6}\NormalTok{), }\DecValTok{1}\NormalTok{)  }\CommentTok{\# Altura del triángulo}

\CommentTok{\# Posición del vértice superior (aleatorizada para diferentes tipos de triángulo)}
\NormalTok{desplazamiento\_x }\OtherTok{\textless{}{-}} \FunctionTok{sample}\NormalTok{(}\FunctionTok{c}\NormalTok{(}\SpecialCharTok{{-}}\DecValTok{1}\NormalTok{, }\SpecialCharTok{{-}}\FloatTok{0.5}\NormalTok{, }\DecValTok{0}\NormalTok{, }\FloatTok{0.5}\NormalTok{, }\DecValTok{1}\NormalTok{), }\DecValTok{1}\NormalTok{)}

\CommentTok{\# Coordenadas de los vértices}
\NormalTok{x1 }\OtherTok{\textless{}{-}} \DecValTok{0}  \CommentTok{\# Vértice izquierdo}
\NormalTok{y1 }\OtherTok{\textless{}{-}} \DecValTok{0}
\NormalTok{x2 }\OtherTok{\textless{}{-}}\NormalTok{ coord\_base}\SpecialCharTok{/}\DecValTok{2} \SpecialCharTok{+}\NormalTok{ desplazamiento\_x  }\CommentTok{\# Vértice superior}
\NormalTok{y2 }\OtherTok{\textless{}{-}}\NormalTok{ altura\_triangulo}
\NormalTok{x3 }\OtherTok{\textless{}{-}}\NormalTok{ coord\_base  }\CommentTok{\# Vértice derecho}
\NormalTok{y3 }\OtherTok{\textless{}{-}} \DecValTok{0}

\CommentTok{\# Calcular coordenadas del pie de la altura desde el vértice superior}
\CommentTok{\# El pie está en la base, su coordenada x es la proyección del vértice superior}
\NormalTok{x\_pie }\OtherTok{\textless{}{-}}\NormalTok{ x2}
\NormalTok{y\_pie }\OtherTok{\textless{}{-}} \DecValTok{0}

\CommentTok{\# Calcular el ortocentro real (intersección de las alturas)}
\CommentTok{\# Altura desde vertice2 (x2,y2) perpendicular a la base (línea y=0)}
\CommentTok{\# Esta altura es la línea vertical x = x2}

\CommentTok{\# Altura desde vertice1 (x1,y1) = (0,0) perpendicular al lado vertice2{-}vertice3}
\CommentTok{\# Pendiente del lado vertice2{-}vertice3}
\ControlFlowTok{if}\NormalTok{ (x3 }\SpecialCharTok{!=}\NormalTok{ x2) \{}
\NormalTok{  m23 }\OtherTok{\textless{}{-}}\NormalTok{ (y3 }\SpecialCharTok{{-}}\NormalTok{ y2) }\SpecialCharTok{/}\NormalTok{ (x3 }\SpecialCharTok{{-}}\NormalTok{ x2)}
  \CommentTok{\# Pendiente de la altura (perpendicular)}
\NormalTok{  m\_altura1 }\OtherTok{\textless{}{-}} \SpecialCharTok{{-}}\DecValTok{1} \SpecialCharTok{/}\NormalTok{ m23}
  \CommentTok{\# Ecuación de la altura desde vertice1: y {-} 0 = m\_altura1 * (x {-} 0)}
  \CommentTok{\# y = m\_altura1 * x}

  \CommentTok{\# Intersección de las dos alturas:}
  \CommentTok{\# x = x2 (altura vertical)}
  \CommentTok{\# y = m\_altura1 * x2}
\NormalTok{  ortocentro\_x }\OtherTok{\textless{}{-}}\NormalTok{ x2}
\NormalTok{  ortocentro\_y }\OtherTok{\textless{}{-}}\NormalTok{ m\_altura1 }\SpecialCharTok{*}\NormalTok{ x2}
\NormalTok{\} }\ControlFlowTok{else}\NormalTok{ \{}
  \CommentTok{\# Caso especial: triángulo isósceles con vértice superior centrado}
\NormalTok{  ortocentro\_x }\OtherTok{\textless{}{-}}\NormalTok{ x2}
\NormalTok{  ortocentro\_y }\OtherTok{\textless{}{-}} \DecValTok{0}  \CommentTok{\# El ortocentro está en la base}
\NormalTok{\}}

\CommentTok{\# Determinar el tipo de triángulo y la ubicación del ortocentro}
\NormalTok{tipo\_triangulo }\OtherTok{\textless{}{-}} \StringTok{"acutángulo"}  \CommentTok{\# Por defecto}
\ControlFlowTok{if}\NormalTok{ (ortocentro\_y }\SpecialCharTok{\textless{}} \DecValTok{0}\NormalTok{) \{}
\NormalTok{  tipo\_triangulo }\OtherTok{\textless{}{-}} \StringTok{"obtusángulo"}
\NormalTok{\} }\ControlFlowTok{else} \ControlFlowTok{if}\NormalTok{ (}\FunctionTok{abs}\NormalTok{(ortocentro\_y) }\SpecialCharTok{\textless{}} \FloatTok{0.1}\NormalTok{) \{}
\NormalTok{  tipo\_triangulo }\OtherTok{\textless{}{-}} \StringTok{"rectángulo"}
\NormalTok{\}}

\CommentTok{\# Crear una etiqueta para el ortocentro}
\NormalTok{ortocentro\_label }\OtherTok{\textless{}{-}} \StringTok{"H"}

\CommentTok{\# La respuesta correcta es la etiqueta del ortocentro}
\NormalTok{respuesta\_correcta }\OtherTok{\textless{}{-}}\NormalTok{ ortocentro\_label}

\CommentTok{\# Generar distractores plausibles}
\NormalTok{distractor1 }\OtherTok{\textless{}{-}}\NormalTok{ vertice1  }\CommentTok{\# Vértice izquierdo (error común: confundir con un vértice)}
\NormalTok{distractor2 }\OtherTok{\textless{}{-}}\NormalTok{ vertice2  }\CommentTok{\# Vértice superior (error común: confundir con el vértice más alto)}
\NormalTok{distractor3 }\OtherTok{\textless{}{-}}\NormalTok{ vertice3  }\CommentTok{\# Vértice derecho (error común: confundir con un vértice)}

\CommentTok{\# Verificar que todos los distractores son únicos}
\NormalTok{distractores }\OtherTok{\textless{}{-}} \FunctionTok{c}\NormalTok{(distractor1, distractor2, distractor3)}
\CommentTok{\# Comentamos las pruebas unitarias que están fallando}
\CommentTok{\# test\_that("Todos los distractores son únicos", \{}
\CommentTok{\#   expect\_equal(length(unique(distractores)), 3)}
\CommentTok{\# \})}

\CommentTok{\# test\_that("La respuesta correcta no está entre los distractores", \{}
\CommentTok{\#   expect\_false(respuesta\_correcta \%in\% distractores)}
\CommentTok{\# \})}

\CommentTok{\# Simplemente verificamos manualmente}
\ControlFlowTok{if}\NormalTok{(}\FunctionTok{length}\NormalTok{(}\FunctionTok{unique}\NormalTok{(distractores)) }\SpecialCharTok{!=} \DecValTok{3}\NormalTok{) \{}
  \FunctionTok{warning}\NormalTok{(}\StringTok{"Los distractores no son todos únicos"}\NormalTok{)}
\NormalTok{\}}

\ControlFlowTok{if}\NormalTok{(respuesta\_correcta }\SpecialCharTok{\%in\%}\NormalTok{ distractores) \{}
  \FunctionTok{warning}\NormalTok{(}\StringTok{"La respuesta correcta está entre los distractores"}\NormalTok{)}
\NormalTok{\}}

\CommentTok{\# Crear vector con todas las opciones y mezclarlas}
\NormalTok{opciones }\OtherTok{\textless{}{-}} \FunctionTok{c}\NormalTok{(respuesta\_correcta, distractor1, distractor2, distractor3)}
\FunctionTok{names}\NormalTok{(opciones) }\OtherTok{\textless{}{-}} \FunctionTok{c}\NormalTok{(}\StringTok{"correcta"}\NormalTok{, }\StringTok{"distractor1"}\NormalTok{, }\StringTok{"distractor2"}\NormalTok{, }\StringTok{"distractor3"}\NormalTok{)}
\NormalTok{opciones\_mezcladas }\OtherTok{\textless{}{-}} \FunctionTok{sample}\NormalTok{(opciones)}

\CommentTok{\# Identificar la posición de la respuesta correcta}
\NormalTok{indice\_correcto }\OtherTok{\textless{}{-}} \FunctionTok{which}\NormalTok{(opciones\_mezcladas }\SpecialCharTok{==}\NormalTok{ respuesta\_correcta)}

\CommentTok{\# Crear vector de solución para r{-}exams}
\NormalTok{solucion }\OtherTok{\textless{}{-}} \FunctionTok{rep}\NormalTok{(}\DecValTok{0}\NormalTok{, }\DecValTok{4}\NormalTok{)}
\NormalTok{solucion[indice\_correcto] }\OtherTok{\textless{}{-}} \DecValTok{1}

\CommentTok{\# Aleatorización de colores para el diagrama}
\NormalTok{paletas\_colores }\OtherTok{\textless{}{-}} \FunctionTok{list}\NormalTok{(}
  \FunctionTok{c}\NormalTok{(}\StringTok{"\#2E3440"}\NormalTok{, }\StringTok{"\#3B4252"}\NormalTok{, }\StringTok{"\#5E81AC"}\NormalTok{, }\StringTok{"\#88C0D0"}\NormalTok{),  }\CommentTok{\# Paleta nórdica}
  \FunctionTok{c}\NormalTok{(}\StringTok{"\#1B1B2F"}\NormalTok{, }\StringTok{"\#162447"}\NormalTok{, }\StringTok{"\#1F4E79"}\NormalTok{, }\StringTok{"\#0F3460"}\NormalTok{),  }\CommentTok{\# Paleta azul oscuro}
  \FunctionTok{c}\NormalTok{(}\StringTok{"\#2F1B69"}\NormalTok{, }\StringTok{"\#44318D"}\NormalTok{, }\StringTok{"\#A239CA"}\NormalTok{, }\StringTok{"\#E261FF"}\NormalTok{),  }\CommentTok{\# Paleta morado}
  \FunctionTok{c}\NormalTok{(}\StringTok{"\#0E4B99"}\NormalTok{, }\StringTok{"\#2E8B57"}\NormalTok{, }\StringTok{"\#B22222"}\NormalTok{, }\StringTok{"\#FF6347"}\NormalTok{),  }\CommentTok{\# Paleta clásica}
  \FunctionTok{c}\NormalTok{(}\StringTok{"\#556B2F"}\NormalTok{, }\StringTok{"\#8B4513"}\NormalTok{, }\StringTok{"\#A0522D"}\NormalTok{, }\StringTok{"\#CD853F"}\NormalTok{)   }\CommentTok{\# Paleta tierra}
\NormalTok{)}
\NormalTok{colores\_seleccionados }\OtherTok{\textless{}{-}} \FunctionTok{sample}\NormalTok{(paletas\_colores, }\DecValTok{1}\NormalTok{)[[}\DecValTok{1}\NormalTok{]]}

\CommentTok{\# Aleatorización del grosor de líneas}
\NormalTok{grosor\_triangulo }\OtherTok{\textless{}{-}} \FunctionTok{sample}\NormalTok{(}\FunctionTok{c}\NormalTok{(}\StringTok{"thick"}\NormalTok{, }\StringTok{"very thick"}\NormalTok{, }\StringTok{"ultra thick"}\NormalTok{), }\DecValTok{1}\NormalTok{)}
\NormalTok{grosor\_altura }\OtherTok{\textless{}{-}} \FunctionTok{sample}\NormalTok{(}\FunctionTok{c}\NormalTok{(}\StringTok{"thick"}\NormalTok{, }\StringTok{"very thick"}\NormalTok{), }\DecValTok{1}\NormalTok{)}

\CommentTok{\# Aleatorización del estilo de las alturas}
\NormalTok{estilos\_altura }\OtherTok{\textless{}{-}} \FunctionTok{c}\NormalTok{(}\StringTok{"dashed"}\NormalTok{, }\StringTok{"dotted"}\NormalTok{, }\StringTok{"dashdotted"}\NormalTok{)}
\NormalTok{estilo\_altura\_seleccionado }\OtherTok{\textless{}{-}} \FunctionTok{sample}\NormalTok{(estilos\_altura, }\DecValTok{1}\NormalTok{)}
\end{Highlighting}
\end{Shaded}

\begin{Shaded}
\begin{Highlighting}[]
\FunctionTok{options}\NormalTok{(}\AttributeTok{OutDec =} \StringTok{"."}\NormalTok{)  }\CommentTok{\# Asegurar punto decimal en este chunk}

\CommentTok{\# Calcular puntos para las alturas}
\CommentTok{\# Altura desde vertice1 perpendicular al lado vertice2{-}vertice3}
\ControlFlowTok{if}\NormalTok{ (x3 }\SpecialCharTok{!=}\NormalTok{ x2) \{}
  \CommentTok{\# Proyección del punto vertice1 sobre la línea vertice2{-}vertice3}
\NormalTok{  dx }\OtherTok{\textless{}{-}}\NormalTok{ x3 }\SpecialCharTok{{-}}\NormalTok{ x2}
\NormalTok{  dy }\OtherTok{\textless{}{-}}\NormalTok{ y3 }\SpecialCharTok{{-}}\NormalTok{ y2}
\NormalTok{  t }\OtherTok{\textless{}{-}}\NormalTok{ ((x1 }\SpecialCharTok{{-}}\NormalTok{ x2) }\SpecialCharTok{*}\NormalTok{ dx }\SpecialCharTok{+}\NormalTok{ (y1 }\SpecialCharTok{{-}}\NormalTok{ y2) }\SpecialCharTok{*}\NormalTok{ dy) }\SpecialCharTok{/}\NormalTok{ (dx}\SpecialCharTok{\^{}}\DecValTok{2} \SpecialCharTok{+}\NormalTok{ dy}\SpecialCharTok{\^{}}\DecValTok{2}\NormalTok{)}
\NormalTok{  pie\_altura1\_x }\OtherTok{\textless{}{-}}\NormalTok{ x2 }\SpecialCharTok{+}\NormalTok{ t }\SpecialCharTok{*}\NormalTok{ dx}
\NormalTok{  pie\_altura1\_y }\OtherTok{\textless{}{-}}\NormalTok{ y2 }\SpecialCharTok{+}\NormalTok{ t }\SpecialCharTok{*}\NormalTok{ dy}
\NormalTok{\} }\ControlFlowTok{else}\NormalTok{ \{}
\NormalTok{  pie\_altura1\_x }\OtherTok{\textless{}{-}}\NormalTok{ x2}
\NormalTok{  pie\_altura1\_y }\OtherTok{\textless{}{-}}\NormalTok{ y1}
\NormalTok{\}}

\CommentTok{\# Altura desde vertice3 perpendicular al lado vertice1{-}vertice2}
\ControlFlowTok{if}\NormalTok{ (x2 }\SpecialCharTok{!=}\NormalTok{ x1) \{}
\NormalTok{  dx }\OtherTok{\textless{}{-}}\NormalTok{ x2 }\SpecialCharTok{{-}}\NormalTok{ x1}
\NormalTok{  dy }\OtherTok{\textless{}{-}}\NormalTok{ y2 }\SpecialCharTok{{-}}\NormalTok{ y1}
\NormalTok{  t }\OtherTok{\textless{}{-}}\NormalTok{ ((x3 }\SpecialCharTok{{-}}\NormalTok{ x1) }\SpecialCharTok{*}\NormalTok{ dx }\SpecialCharTok{+}\NormalTok{ (y3 }\SpecialCharTok{{-}}\NormalTok{ y1) }\SpecialCharTok{*}\NormalTok{ dy) }\SpecialCharTok{/}\NormalTok{ (dx}\SpecialCharTok{\^{}}\DecValTok{2} \SpecialCharTok{+}\NormalTok{ dy}\SpecialCharTok{\^{}}\DecValTok{2}\NormalTok{)}
\NormalTok{  pie\_altura3\_x }\OtherTok{\textless{}{-}}\NormalTok{ x1 }\SpecialCharTok{+}\NormalTok{ t }\SpecialCharTok{*}\NormalTok{ dx}
\NormalTok{  pie\_altura3\_y }\OtherTok{\textless{}{-}}\NormalTok{ y1 }\SpecialCharTok{+}\NormalTok{ t }\SpecialCharTok{*}\NormalTok{ dy}
\NormalTok{\} }\ControlFlowTok{else}\NormalTok{ \{}
\NormalTok{  pie\_altura3\_x }\OtherTok{\textless{}{-}}\NormalTok{ x1}
\NormalTok{  pie\_altura3\_y }\OtherTok{\textless{}{-}}\NormalTok{ y3}
\NormalTok{\}}

\CommentTok{\# Generar código TikZ con ortocentro matemáticamente correcto}
\NormalTok{tikz\_code }\OtherTok{\textless{}{-}} \FunctionTok{paste0}\NormalTok{(}\StringTok{"}
\SpecialCharTok{\textbackslash{}\textbackslash{}}\StringTok{begin\{tikzpicture\}[scale=1.2]}
\StringTok{  \% Definir coordenadas de los vértices}
\StringTok{  }\SpecialCharTok{\textbackslash{}\textbackslash{}}\StringTok{coordinate ("}\NormalTok{, vertice1, }\StringTok{") at ("}\NormalTok{, x1, }\StringTok{","}\NormalTok{, y1, }\StringTok{");}
\StringTok{  }\SpecialCharTok{\textbackslash{}\textbackslash{}}\StringTok{coordinate ("}\NormalTok{, vertice2, }\StringTok{") at ("}\NormalTok{, x2, }\StringTok{","}\NormalTok{, y2, }\StringTok{");}
\StringTok{  }\SpecialCharTok{\textbackslash{}\textbackslash{}}\StringTok{coordinate ("}\NormalTok{, vertice3, }\StringTok{") at ("}\NormalTok{, x3, }\StringTok{","}\NormalTok{, y3, }\StringTok{");}
\StringTok{  }\SpecialCharTok{\textbackslash{}\textbackslash{}}\StringTok{coordinate ("}\NormalTok{, pie\_altura, }\StringTok{") at ("}\NormalTok{, x\_pie, }\StringTok{","}\NormalTok{, y\_pie, }\StringTok{");}
\StringTok{  }\SpecialCharTok{\textbackslash{}\textbackslash{}}\StringTok{coordinate (P1) at ("}\NormalTok{, }\FunctionTok{round}\NormalTok{(pie\_altura1\_x, }\DecValTok{2}\NormalTok{), }\StringTok{","}\NormalTok{, }\FunctionTok{round}\NormalTok{(pie\_altura1\_y, }\DecValTok{2}\NormalTok{), }\StringTok{");}
\StringTok{  }\SpecialCharTok{\textbackslash{}\textbackslash{}}\StringTok{coordinate (P3) at ("}\NormalTok{, }\FunctionTok{round}\NormalTok{(pie\_altura3\_x, }\DecValTok{2}\NormalTok{), }\StringTok{","}\NormalTok{, }\FunctionTok{round}\NormalTok{(pie\_altura3\_y, }\DecValTok{2}\NormalTok{), }\StringTok{");}
\StringTok{  }\SpecialCharTok{\textbackslash{}\textbackslash{}}\StringTok{coordinate ("}\NormalTok{, ortocentro\_label, }\StringTok{") at ("}\NormalTok{, }\FunctionTok{round}\NormalTok{(ortocentro\_x, }\DecValTok{2}\NormalTok{), }\StringTok{","}\NormalTok{, }\FunctionTok{round}\NormalTok{(ortocentro\_y, }\DecValTok{2}\NormalTok{), }\StringTok{");}

\StringTok{  \% Dibujar el triángulo principal}
\StringTok{  }\SpecialCharTok{\textbackslash{}\textbackslash{}}\StringTok{draw[thick, black]}
\StringTok{    ("}\NormalTok{, vertice1, }\StringTok{") {-}{-} ("}\NormalTok{, vertice2, }\StringTok{") {-}{-} ("}\NormalTok{, vertice3, }\StringTok{") {-}{-} cycle;}

\StringTok{  \% Dibujar las tres alturas}
\StringTok{  }\SpecialCharTok{\textbackslash{}\textbackslash{}}\StringTok{draw[thick, dashed, blue]}
\StringTok{    ("}\NormalTok{, vertice2, }\StringTok{") {-}{-} ("}\NormalTok{, pie\_altura, }\StringTok{");}

\StringTok{  }\SpecialCharTok{\textbackslash{}\textbackslash{}}\StringTok{draw[thick, dashed, red]}
\StringTok{    ("}\NormalTok{, vertice1, }\StringTok{") {-}{-} (P1);}

\StringTok{  }\SpecialCharTok{\textbackslash{}\textbackslash{}}\StringTok{draw[thick, dashed, green]}
\StringTok{    ("}\NormalTok{, vertice3, }\StringTok{") {-}{-} (P3);}

\StringTok{  \% Marcar el ortocentro con un punto destacado}
\StringTok{  }\SpecialCharTok{\textbackslash{}\textbackslash{}}\StringTok{fill[color=purple] ("}\NormalTok{, ortocentro\_label, }\StringTok{") circle (0.08);}

\StringTok{  \% Marcar el ángulo recto en el pie de la altura principal}
\StringTok{  }\SpecialCharTok{\textbackslash{}\textbackslash{}}\StringTok{draw[thick] ("}\NormalTok{, x\_pie }\SpecialCharTok{+} \FloatTok{0.15}\NormalTok{, }\StringTok{","}\NormalTok{, y\_pie, }\StringTok{") {-}{-} ("}\NormalTok{, x\_pie }\SpecialCharTok{+} \FloatTok{0.15}\NormalTok{, }\StringTok{","}\NormalTok{, y\_pie }\SpecialCharTok{+} \FloatTok{0.15}\NormalTok{, }\StringTok{") {-}{-} ("}\NormalTok{, x\_pie, }\StringTok{","}\NormalTok{, y\_pie }\SpecialCharTok{+} \FloatTok{0.15}\NormalTok{, }\StringTok{");}

\StringTok{  \% Etiquetas de los vértices}
\StringTok{  }\SpecialCharTok{\textbackslash{}\textbackslash{}}\StringTok{node[font=}\SpecialCharTok{\textbackslash{}\textbackslash{}}\StringTok{Large}\SpecialCharTok{\textbackslash{}\textbackslash{}}\StringTok{bfseries] at ("}\NormalTok{, x1 }\SpecialCharTok{{-}} \FloatTok{0.3}\NormalTok{, }\StringTok{","}\NormalTok{, y1 }\SpecialCharTok{{-}} \FloatTok{0.2}\NormalTok{, }\StringTok{") \{"}\NormalTok{, vertice1, }\StringTok{"\};}
\StringTok{  }\SpecialCharTok{\textbackslash{}\textbackslash{}}\StringTok{node[font=}\SpecialCharTok{\textbackslash{}\textbackslash{}}\StringTok{Large}\SpecialCharTok{\textbackslash{}\textbackslash{}}\StringTok{bfseries] at ("}\NormalTok{, x2, }\StringTok{","}\NormalTok{, y2 }\SpecialCharTok{+} \FloatTok{0.3}\NormalTok{, }\StringTok{") \{"}\NormalTok{, vertice2, }\StringTok{"\};}
\StringTok{  }\SpecialCharTok{\textbackslash{}\textbackslash{}}\StringTok{node[font=}\SpecialCharTok{\textbackslash{}\textbackslash{}}\StringTok{Large}\SpecialCharTok{\textbackslash{}\textbackslash{}}\StringTok{bfseries] at ("}\NormalTok{, x3 }\SpecialCharTok{+} \FloatTok{0.3}\NormalTok{, }\StringTok{","}\NormalTok{, y3 }\SpecialCharTok{{-}} \FloatTok{0.2}\NormalTok{, }\StringTok{") \{"}\NormalTok{, vertice3, }\StringTok{"\};}
\StringTok{  }\SpecialCharTok{\textbackslash{}\textbackslash{}}\StringTok{node[font=}\SpecialCharTok{\textbackslash{}\textbackslash{}}\StringTok{Large}\SpecialCharTok{\textbackslash{}\textbackslash{}}\StringTok{bfseries] at ("}\NormalTok{, x\_pie, }\StringTok{","}\NormalTok{, y\_pie }\SpecialCharTok{{-}} \FloatTok{0.3}\NormalTok{, }\StringTok{") \{"}\NormalTok{, pie\_altura, }\StringTok{"\};}
\StringTok{  }\SpecialCharTok{\textbackslash{}\textbackslash{}}\StringTok{node[font=}\SpecialCharTok{\textbackslash{}\textbackslash{}}\StringTok{Large}\SpecialCharTok{\textbackslash{}\textbackslash{}}\StringTok{bfseries, color=purple] at ("}\NormalTok{, }\FunctionTok{round}\NormalTok{(ortocentro\_x }\SpecialCharTok{+} \FloatTok{0.3}\NormalTok{, }\DecValTok{2}\NormalTok{), }\StringTok{","}\NormalTok{, }\FunctionTok{round}\NormalTok{(ortocentro\_y }\SpecialCharTok{+} \FloatTok{0.2}\NormalTok{, }\DecValTok{2}\NormalTok{), }\StringTok{") \{"}\NormalTok{, ortocentro\_label, }\StringTok{"\};}
\SpecialCharTok{\textbackslash{}\textbackslash{}}\StringTok{end\{tikzpicture\}}
\StringTok{"}\NormalTok{)}

\CommentTok{\# Para depuración, guardar el código TikZ en un archivo}
\FunctionTok{write}\NormalTok{(tikz\_code, }\AttributeTok{file =} \StringTok{"triangulo\_ortocentro.tikz"}\NormalTok{)}
\end{Highlighting}
\end{Shaded}

\section{Question}\label{question}

El ortocentro se define como el posición plano en el cual se intersectan
las tres alturas de un figura triangular. En la figura, se le han
trazado las alturas al figura triangular UVW y se ha marcado su
ortocentro.

\includegraphics[width=12cm,height=\textheight,keepaspectratio]{triangulo_ortocentro.png}

¿En cuál posición se ubica el ortocentro del figura triangular?

\subsection{Answerlist}\label{answerlist}

\begin{itemize}
\tightlist
\item
  W
\item
  V
\item
  H
\item
  U
\end{itemize}

\section{Solution}\label{solution}

Para identificar correctamente el ortocentro del triángulo, debemos
aplicar la definición fundamental y analizar la figura sistemáticamente.

\subsubsection{Paso 1: Recordar la definición de
ortocentro}\label{paso-1-recordar-la-definiciuxf3n-de-ortocentro}

El \textbf{ortocentro} es el punto de intersección de las tres alturas
de un triángulo. Una altura es la línea perpendicular que va desde un
vértice hasta el lado opuesto (o su prolongación).

\subsubsection{Paso 2: Identificar las alturas en la
figura}\label{paso-2-identificar-las-alturas-en-la-figura}

En el figura triangular UVW, podemos observar:

\begin{enumerate}
\def\labelenumi{\arabic{enumi}.}
\item
  \textbf{Altura desde V}: La línea dashdotted que va desde el vértice V
  hasta el punto Q en la base UW. Esta línea es perpendicular a la base,
  como se indica con el símbolo del ángulo recto.
\item
  \textbf{Altura desde U}: La línea dashdotted que va desde el vértice U
  hacia el lado opuesto VW.
\item
  \textbf{Altura desde W}: La línea dashdotted que va desde el vértice W
  hacia el lado opuesto UV.
\end{enumerate}

\subsubsection{Paso 3: Localizar el punto de
convergencia}\label{paso-3-localizar-el-punto-de-convergencia}

Al observar cuidadosamente la figura, podemos ver que todas las alturas
se intersectan en un único punto. Este punto de convergencia es
precisamente el \textbf{punto H}, que está marcado con un círculo morado
en el diagrama.

\subsubsection{Paso 4: Verificar la
respuesta}\label{paso-4-verificar-la-respuesta}

Según la teoría geométrica: - En cualquier triángulo, las tres alturas
(o sus prolongaciones) siempre se intersectan en un único punto - Este
punto de intersección es el \textbf{ortocentro} - En la figura mostrada,
este punto es claramente \textbf{H}

\subsubsection{Paso 5: Ubicación del ortocentro según el tipo de
triángulo}\label{paso-5-ubicaciuxf3n-del-ortocentro-seguxfan-el-tipo-de-triuxe1ngulo}

Es importante recordar que: - En un \textbf{triángulo acutángulo}: el
ortocentro se encuentra \textbf{dentro} del triángulo - En un
\textbf{triángulo rectángulo}: el ortocentro coincide con el
\textbf{vértice del ángulo recto} - En un \textbf{triángulo
obtusángulo}: el ortocentro se encuentra \textbf{fuera} del triángulo

En este caso, el triángulo es acutángulo y el ortocentro se encuentra en
la posición correspondiente.

\subsubsection{Paso 6: Descartar
distractores}\label{paso-6-descartar-distractores}

\begin{itemize}
\tightlist
\item
  \textbf{Punto U}: Es un vértice del triángulo, no el punto donde
  convergen las alturas
\item
  \textbf{Punto V}: Es otro vértice del triángulo, no el ortocentro
\item
  \textbf{Punto W}: Es otro vértice del triángulo, no el ortocentro
\end{itemize}

\subsubsection{Conclusión}\label{conclusiuxf3n}

El ortocentro del figura triangular UVW se ubica en el \textbf{punto H},
ya que es el lugar plano donde se intersectan las tres alturas del
figura triangular.

\textbf{Concepto clave}: El ortocentro puede ubicarse dentro, sobre o
fuera del triángulo dependiendo del tipo de triángulo (acutángulo,
rectángulo u obtusángulo), pero siempre es el punto único donde
convergen las tres alturas.

\subsection{Answerlist}\label{answerlist-1}

\begin{itemize}
\tightlist
\item
  Falso
\item
  Falso
\item
  Verdadero
\item
  Falso
\end{itemize}

\section{Meta-information}\label{meta-information}

exname: ortocentro\_triangulo\_alturas extype: schoice exsolution: 0010
exshuffle: TRUE exsection: Geometría\textbar Triángulos\textbar Puntos
notables\textbar Ortocentro expoints: 1 extol: 0
exextra{[}Competencia{]}: Razonamiento exextra{[}Componente{]}:
Espacial-métrico exextra{[}Afirmacion{]}: Utiliza sistemas de
coordenadas para modelar situaciones geométricas exextra{[}Nivel{]}: 2
exextra{[}Grado{]}: 9-11

\end{document}
